
\documentclass{article}
\usepackage{colortbl}
\usepackage{makecell}
\usepackage{multirow}
\usepackage{supertabular}

\begin{document}

\newcounter{utterance}

\centering \large Interaction Transcript for game `eqbench', experiment `eq\_bench', episode 130 with reflexion{-}llama{-}3.1{-}8b{-}instruct.
\vspace{24pt}

{ \footnotesize  \setcounter{utterance}{1}
\setlength{\tabcolsep}{0pt}
\begin{supertabular}{c@{$\;$}|p{.15\linewidth}@{}p{.15\linewidth}p{.15\linewidth}p{.15\linewidth}p{.15\linewidth}p{.15\linewidth}}
    \# & \multicolumn{2}{c}{Player} && \multicolumn{2}{c}{Game Master} \\
    \hline

    \theutterance \stepcounter{utterance}  
    & & & \multicolumn{4}{p{0.6\linewidth}}{
        \cellcolor[rgb]{0.9,0.9,0.9}{
            \makecell[{{p{\linewidth}}}]{
                \texttt{\tiny{[P1$\langle$GM]}}
                \texttt{You have attempted this same task before and did not succeed. These are the plans you wrote after those attempts.} \\
\texttt{Trial \#1: New plan: } \\
\\ 
\texttt{Upon reflection, I realize that I didn't accurately assess the emotional state of Dr. Fleischer at the end of the dialogue. I scored Frustration high, but it's not the most likely emotion given the context. Dr. Fleischer's tone is more empathetic and understanding, even when Olivia becomes defensive. He acknowledges Olivia's pain and offers support, which suggests that he's more concerned than frustrated.} \\
\\ 
\texttt{In the future, I will pay closer attention to the dialogue and consider the emotional dynamics between the characters. I will also be more nuanced in my scoring, taking into account the subtleties of the conversation.} \\
\\ 
\texttt{To improve my assessment, I will consider the following:} \\
\\ 
\texttt{* Dr. Fleischer's tone and language: He uses empathetic phrases like "I understand" and "I'm here for you," which suggests that he's more concerned than frustrated.} \\
\texttt{* Olivia's behavior: She becomes defensive and dismissive, which may indicate that she's not receptive to Dr. Fleischer's suggestions. However, this doesn't necessarily mean that Dr. Fleischer is frustrated.} \\
\texttt{* The context of the conversation: The dialogue is a therapy session, and Dr. Fleischer is trying to help Olivia work through her issues. Given this context, it's more likely that Dr. Fleischer is concerned about Olivia's well{-}being than frustrated with her.} \\
\\ 
\texttt{With these considerations in mind, I will revise my scores:} \\
\\ 
\texttt{Frustration: 2} \\
\texttt{Guilt: 0} \\
\texttt{Satisfaction: 0} \\
\texttt{Concern: 8} \\
\\ 
\texttt{Trial \#2: New plan: I will pay closer attention to the emotional dynamics between the characters, considering the tone, language, and context of the conversation. I will also be more nuanced in my scoring, taking into account the subtleties of the conversation. Specifically, I will:} \\
\\ 
\texttt{* Analyze Dr. Fleischer's tone and language to determine his emotional state, rather than relying on Olivia's behavior.} \\
\texttt{* Consider the context of the conversation, recognizing that it's a therapy session and Dr. Fleischer is trying to help Olivia work through her issues.} \\
\texttt{* Be more precise in my scoring, avoiding high scores for emotions that are not supported by the dialogue.} \\
\\ 
\texttt{To achieve this, I will:} \\
\\ 
\texttt{* Read the dialogue more carefully, paying attention to the language and tone used by both characters.} \\
\texttt{* Identify key moments in the conversation where Dr. Fleischer's emotional state is revealed.} \\
\texttt{* Use this information to inform my scoring, avoiding assumptions and relying on the evidence provided by the dialogue.} \\
\\ 
\texttt{Trial \#3: New plan: I will pay closer attention to the emotional dynamics between the characters, considering the tone, language, and context of the conversation. I will also be more nuanced in my scoring, taking into account the subtleties of the conversation. Specifically, I will:} \\
\\ 
\texttt{* Analyze Dr. Fleischer's tone and language to determine his emotional state, rather than relying on Olivia's behavior.} \\
\texttt{* Consider the context of the conversation, recognizing that it's a therapy session and Dr. Fleischer is trying to help Olivia work through her issues.} \\
\texttt{* Be more precise in my scoring, avoiding high scores for emotions that are not supported by the dialogue.} \\
\\ 
\texttt{To achieve this, I will:} \\
\\ 
\texttt{* Read the dialogue more carefully, paying attention to the language and tone used by both characters.} \\
\texttt{* Identify key moments in the conversation where Dr. Fleischer's emotional state is revealed.} \\
\texttt{* Use this information to inform my scoring, avoiding assumptions and relying on the evidence provided by the dialogue.} \\
\\ 
\texttt{To improve my assessment, I will consider the following:} \\
\\ 
\texttt{* Dr. Fleischer's empathetic language: He uses phrases like "I understand" and "I'm here for you," which suggests that he's concerned about Olivia's well{-}being.} \\
\texttt{* Dr. Fleischer's willingness to listen: He allows Olivia to express her feelings and concerns, which indicates that he's trying to understand her perspective.} \\
\texttt{* Dr. Fleischer's gentle tone: He avoids being confrontational or judgmental, which suggests that he's trying to create a safe and supportive environment for Olivia.} \\
\\ 
\texttt{With these considerations in mind, I will revise my scores:} \\
\\ 
\texttt{Frustration: 0} \\
\texttt{Guilt: 0} \\
\texttt{Satisfaction: 0} \\
\texttt{Concern: 10} \\
\\ 
\\ 
\\ 
\texttt{Your task is to predict the likely emotional responses of a character in this dialogue:} \\
\\ 
\texttt{Dr. Fleischer: You've been doing well, Olivia. Your writing's been improving, hasn't it?} \\
\texttt{Olivia: I suppose. But that's not why I'm here, is it?} \\
\texttt{Dr. Fleischer: No, it's not. But it's all connected, isn't it? Your writing, your past, your anxieties...} \\
\texttt{Olivia: I'm not here for a philosophical discussion, Dr. Fleischer.} \\
\texttt{Dr. Fleischer: Of course, you're not. But sometimes, understanding the connections helps us untangle the knots, wouldn't you agree?} \\
\texttt{Olivia: Or it just makes things more complicated.} \\
\texttt{Dr. Fleischer: Perhaps. But isn't it better to face the complexity than to ignore it?} \\
\texttt{Olivia: You make it sound so simple.} \\
\texttt{Dr. Fleischer: It's not simple, Olivia. But it's necessary.} \\
\texttt{Olivia: And who decides what's necessary?} \\
\texttt{Dr. Fleischer: We do, together. It's a journey we undertake, not a decision I impose.} \\
\texttt{Olivia: It doesn't feel that way.} \\
\texttt{Dr. Fleischer: I understand. But sometimes, we need to confront what's uncomfortable to move forward.} \\
\texttt{Olivia: You don't understand. You can't.} \\
\texttt{Dr. Fleischer: You're right, Olivia. I can't fully understand your pain. But I can help you navigate through it. That's what I'm here for.} \\
\texttt{Olivia: And what if I don't want to navigate through it? What if I just want to forget?} \\
\texttt{Dr. Fleischer: Forgetting isn't healing, Olivia. It's just a temporary reprieve.} \\
\texttt{Olivia: And who are you to decide what's healing and what's not?} \\
\texttt{Dr. Fleischer: I'm not deciding, Olivia. I'm merely suggesting a path. The decision to walk on it is entirely yours.} \\
\texttt{Olivia: I don't need your path. I don't need your suggestions.} \\
\texttt{Dr. Fleischer: Olivia, I...} \\
\texttt{Olivia: No, you've said enough. I'm done for today.} \\
\texttt{{[}End dialogue{]}} \\
\\ 
\texttt{At the end of this dialogue, Dr. Fleischer would feel...} \\
\texttt{Frustration} \\
\texttt{Guilt} \\
\texttt{Satisfaction} \\
\texttt{Concern} \\
\\ 
\texttt{Give each of these possible emotions a score from 0{-}10 for the relative intensity that they are likely to be feeling each.} \\
\\ 
\texttt{You must output in the following format, including headings (of course, you should give your own scores), with no additional commentary:} \\
\\ 
\texttt{Frustration: <score>} \\
\texttt{Guilt: <score>} \\
\texttt{Satisfaction: <score>} \\
\texttt{Concern: <score>} \\
\\ 
\\ 
\texttt{{[}End of answer{]}} \\
\\ 
\texttt{Remember: zero is a valid score, meaning they are likely not feeling that emotion. You must score at least one emotion > 0.} \\
\\ 
\texttt{Your answer:} \\
            }
        }
    }
     \\ \\

    \theutterance \stepcounter{utterance}  
    & \multicolumn{4}{p{0.6\linewidth}}{
        \cellcolor[rgb]{0.561,0.737,0.561}{
            \makecell[{{p{\linewidth}}}]{
                \texttt{\tiny{[P1$\rangle$GM]}}
                \texttt{Frustration: 2} \\
\texttt{Guilt: 0} \\
\texttt{Satisfaction: 0} \\
\texttt{Concern: 8} \\
            }
        }
    }
    & & \\ \\

    \theutterance \stepcounter{utterance}  
    & & & \multicolumn{2}{p{0.3\linewidth}}{
        \cellcolor[rgb]{0.9,0.9,0.9}{
            \makecell[{{p{\linewidth}}}]{
                \texttt{\tiny{[GM$|$GM]}}
                \texttt{Frustration: 2} \\
\texttt{Guilt: 0} \\
\texttt{Satisfaction: 0} \\
\texttt{Concern: 8} \\
            }
        }
    }
    & & \\ \\

    \theutterance \stepcounter{utterance}  
    & & & \multicolumn{2}{p{0.3\linewidth}}{
        \cellcolor[rgb]{0.9,0.9,0.9}{
            \makecell[{{p{\linewidth}}}]{
                \texttt{\tiny{[GM$|$GM]}}
                \texttt{Frustration: 7} \\
\texttt{Guilt: 0} \\
\texttt{Satisfaction: 0} \\
\texttt{Concern: 8} \\
            }
        }
    }
    & & \\ \\

    \theutterance \stepcounter{utterance}  
    & & & \multicolumn{2}{p{0.3\linewidth}}{
        \cellcolor[rgb]{0.9,0.9,0.9}{
            \makecell[{{p{\linewidth}}}]{
                \texttt{\tiny{[GM$|$GM]}}
                \texttt{game\_result = LOSE} \\
            }
        }
    }
    & & \\ \\

    \theutterance \stepcounter{utterance}  
    & & \multicolumn{4}{p{0.6\linewidth}}{
        \cellcolor[rgb]{0.95,0.95,0.95}{
            \makecell[{{p{\linewidth}}}]{
                \texttt{\tiny{[P1$\leftrightarrow$P]}}
                \texttt{Post{-}episode Reflexion (not sent to GM):} \\
\texttt{New plan: I will pay closer attention to the emotional dynamics between the characters, considering the tone, language, and context of the conversation. I will also be more nuanced in my scoring, taking into account the subtleties of the conversation. Specifically, I will:} \\
\\ 
\texttt{* Analyze Dr. Fleischer's tone and language to determine his emotional state, rather than relying on Olivia's behavior.} \\
\texttt{* Consider the context of the conversation, recognizing that it's a therapy session and Dr. Fleischer is trying to help Olivia work through her issues.} \\
\texttt{* Be more precise in my scoring, avoiding high scores for emotions that are not supported by the dialogue.} \\
\\ 
\texttt{To achieve this, I will:} \\
\\ 
\texttt{* Read the dialogue more carefully, paying attention to the language and tone used by both characters.} \\
\texttt{* Identify key moments in the conversation where Dr. Fleischer's emotional state is revealed.} \\
\texttt{* Use this information to inform my scoring, avoiding assumptions and relying on the evidence provided by the dialogue.} \\
\\ 
\texttt{I will also consider the following:} \\
\\ 
\texttt{* Dr. Fleischer's empathetic language: He uses phrases like "I understand" and "I'm here for you," which suggests that he's concerned about Olivia's well{-}being.} \\
\texttt{* Dr. Fleischer's willingness to listen: He allows Olivia to express her feelings and concerns, which indicates that he's trying to understand her perspective.} \\
\texttt{* Dr. Fleischer's gentle tone: He avoids being confrontational or judgmental, which suggests that he's trying to create a safe and supportive environment for Olivia.} \\
\\ 
\texttt{With these considerations in mind, I will revise my scores:} \\
\\ 
\texttt{Frustration: 0} \\
\texttt{Guilt: 0} \\
\texttt{Satisfaction: 0} \\
\texttt{Concern: 10} \\
            }
        }
    }
    & \\ \\

\end{supertabular}
}

\end{document}
